\SourceLeaf{63}
mit $\left(1-\dfrac{r}{\nu}\right)$. Die Multiplikation mit $\dfrac{r}{\nu}$ kann man dadurch ausführen, dass
man das vorherige Glied nach $n$ differenziert, erneut mit $n$ multipliziert und durch $\nu\,\mathrm{d}n$ dividiert.
Damit ergibt sich die folgende Beziehung für das Anfangsglied von $B^{(r)}$
\[
= A^{(r)}\text{ term. init.} - \frac{n\,{}^{*}\!\mathrm{d}A^{(r)}\text{ term. init.}}{\nu\,\mathrm{d}n}.
\]
Genau dieselbe Beziehung gilt für die allgemeinen Glieder. Für den Koeffizienten
$A^{(r)}$ lautet das Glied mit Index $h^{\mathrm{tus}}$
\[
= \frac{1}{2^{r+2h-1}}\,{}^{r+2h}\mathfrak{B}^{(r+h)} . {}^{\nu}\mathfrak{B}^{(r+2h)}\,n^{r+2h}.
\]
Ändert man also $\nu$ in $\nu-1$, so erhält man für $B^{(r)}$ das Glied mit Index $h^{\mathrm{tus}}$
\[
= \frac{1}{2^{r+2h-1}}\,{}^{r+2h}\mathfrak{B}^{(r+h)} . {}^{\nu-1}\mathfrak{B}^{(r+2h)}\,n^{r+2h},
\]
das aus dem vorherigen Ausdruck durch bloßes Hinzufügen des Faktors
\[
\frac{\nu-(r+2h)}{\nu} = 1 - \frac{r+2h}{\nu}
\]
unmittelbar hervorgeht. Das Glied von
$A^{(r)}$ mit Index $h^{\mathrm{tum}}$ mit $\dfrac{r+2h}{\nu}$ zu multiplizieren ist gleichbedeutend damit, es nach
$n$ zu differenzieren, durch $\nu\,\mathrm{d}n$ zu dividieren und erneut mit $n$ zu multiplizieren. Damit erhalten wir
\[
B^{(r)}\text{ term. }h = A^{(r)}\text{ term. }h - \frac{n\,{}^{*}\!\mathrm{d}A^{(r)}\text{ term. }h}{\nu\,\mathrm{d}n},
\]
und da diese Beziehung ausnahmslos für alle Glieder gilt, ergibt sich die Gleichung
\[
B^{(r)} = A^{(r)} - \frac{n\,{}^{*}\!\mathrm{d}A^{(r)}}{\nu\,\mathrm{d}n}.
\]
\par\noindent\hspace*{\parindent}$\varepsilon$) Durch diese Rekursionsformel wird also $(1+n\cos\varphi)^{\nu-1}$ auf $(1+n\cos\varphi)^{\nu}$
so zurückgeführt, dass sich die Koeffizienten des ersten Ausdrucks durch Differentiation sehr leicht
aus den im zweiten Ausdruck vorkommenden ableiten lassen. Da uns die Koeffizienten von $(1+n\cos\varphi)^{-1}$
in endlicher Form vorliegen, können wir schrittweise zu negativen Exponenten größeren Betrags übergehen
und somit ganzzahlige negative Potenzen allgemein durch endliche Koeffizienten darstellen.

Die zuvor für diese Koeffizienten gefundenen unendlichen Reihen sind dennoch von Nutzen.
Sie kommen zur Anwendung, sobald $\nu$ eine gebrochene Zahl ist. In diesem Fall darf man nach dem in diesen
Untersuchungen besonders gewichtigen Urteil \textsc{Euler}s in keiner Weise erwarten, endliche Formeln
